\documentclass[12pt,a4paper]{article}

% ============================================================================
% Encodage et langue
% ============================================================================
\usepackage[utf8]{inputenc}
\usepackage[T1]{fontenc}
\usepackage[french]{babel}
\DeclareUnicodeCharacter{202F}{\,}

% ============================================================================
% Mise en page
% ============================================================================
\usepackage[a4paper,margin=2.25cm]{geometry}
\usepackage{setspace}
\onehalfspacing
\usepackage{microtype}

% ============================================================================
% Maths
% ============================================================================
\usepackage{amsmath,amssymb,amsthm}
\usepackage{newtxtext,newtxmath}

% ============================================================================
% Figures, couleurs et graphes
% ============================================================================
\usepackage{graphicx}
\usepackage{xcolor}
\usepackage{tikz}
\usetikzlibrary{arrows.meta,positioning,calc,babel,shapes.geometric}
\usepackage{pgfplots}
\pgfplotsset{compat=1.18}

% ============================================================================
% Tableaux
% ============================================================================
\usepackage{booktabs}
\usepackage{array}
\usepackage{multirow}

% ============================================================================
% Divers
% ============================================================================
\usepackage{caption}
\usepackage{subcaption}
\usepackage{float}
\usepackage{listings}
\usepackage{hyperref}
\usepackage{fancyhdr}
\usepackage{lastpage}
\usepackage{tocloft}
\usepackage{titlesec}
\usepackage{enumitem}

% ============================================================================
% Liens : table des matières en noir
% ============================================================================
\hypersetup{
    colorlinks=true,
    linkcolor=black,
    urlcolor=blue!55!black,
    citecolor=black
}

% Toutes les images doivent être dans images/
\graphicspath{{images/}}

% ============================================================================
% Réglages généraux
% ============================================================================
\setcounter{tocdepth}{2}
\setlength{\parindent}{1.15cm}
\setlength{\parskip}{0pt}

\fancypagestyle{rapportstyle}{%
    \fancyhf{}
    \fancyhead[L]{\small Simulation d'un Google Bombing}
    \fancyhead[R]{\small M1 Informatique}
    \renewcommand{\headrulewidth}{0.4pt}
    \fancyfoot[C]{\thepage/\pageref{LastPage}}
}

\renewcommand{\contentsname}{Table des matières}
\renewcommand{\cfttoctitlefont}{\Huge\bfseries}
\renewcommand{\cftsecfont}{\bfseries}
\renewcommand{\cftsecpagefont}{\bfseries}
\renewcommand{\cftsecleader}{\hfill}
\renewcommand{\cftsubsecleader}{\cftdotfill{\cftdotsep}}
\setlength{\cftbeforesecskip}{0.65em}
\setlength{\cftbeforesubsecskip}{0.15em}
\setlength{\cftsecnumwidth}{1.55em}
\setlength{\cftsubsecindent}{1.55em}
\setlength{\cftsubsecnumwidth}{2.65em}

\titleformat{\section}{\normalfont\Large\bfseries}{\thesection}{1em}{}
\titleformat{\subsection}{\normalfont\large\bfseries}{\thesubsection}{1em}{}
\titlespacing*{\section}{0pt}{2.2em}{1.0em}
\titlespacing*{\subsection}{0pt}{1.4em}{0.65em}

% ============================================================================
% Commandes perso
% ============================================================================
\newcommand{\PR}{\mathrm{PR}}
\newcommand{\vpi}{\boldsymbol{\pi}}
\newcommand{\G}{\mathbf{G}}
\newcommand{\Hmat}{\mathbf{H}}

\theoremstyle{definition}
\newtheorem{definition}{Définition}
\newtheorem{remarque}{Remarque}[section]

% ============================================================================
% Logo UVSQ
% ============================================================================
\newcommand{\LogoUVSQ}{%
  \IfFileExists{images/logo-uvsq-2020-rvb.png}{%
    \includegraphics[width=2.65cm]{images/logo-uvsq-2020-rvb.png}%
  }{%
    \fbox{\parbox[c][1.75cm][c]{2.65cm}{\centering\bfseries UVSQ\\Paris-Saclay}}%
  }%
}

% ============================================================================
% Listings
% ============================================================================
\lstset{
    basicstyle=\ttfamily\footnotesize,
    keywordstyle=\color{blue!55!black}\bfseries,
    commentstyle=\color{green!35!black},
    stringstyle=\color{orange!75!black},
    showstringspaces=false,
    breaklines=true,
    frame=single,
    rulecolor=\color{black!25},
    tabsize=2,
    columns=fullflexible
}

% ============================================================================
% Polices de la page de garde
% ============================================================================
\newcommand{\HeaderFont}{\fontsize{17}{21}\selectfont}
\newcommand{\MainTitleFont}{\fontsize{25}{30}\selectfont}
\newcommand{\SubTitleFont}{\fontsize{18}{22}\selectfont}
\newcommand{\BlockTitleFont}{\fontsize{17}{21}\selectfont}
\newcommand{\FormationFont}{\fontsize{15}{18}\selectfont}
\newcommand{\AuthorTitleFont}{\fontsize{16}{19}\selectfont}
\newcommand{\AuthorFont}{\fontsize{13}{17}\selectfont}
\newcommand{\ResponsibleTitleFont}{\fontsize{15}{18}\selectfont}
\newcommand{\ResponsibleFont}{\fontsize{13}{17}\selectfont}
\newcommand{\YearFont}{\fontsize{13}{16}\selectfont}

\begin{document}

% ============================================================================
% Page de garde
% ============================================================================
\begin{titlepage}
\newgeometry{margin=0cm}
\thispagestyle{empty}
\begin{tikzpicture}[remember picture,overlay]

\node[anchor=center] at ([xshift=2.75cm,yshift=-2.80cm]current page.north west)
{\LogoUVSQ};

\node[anchor=center,align=center] at ([yshift=-2.80cm]current page.north)
{\HeaderFont\bfseries\itshape Université de Versailles Saint-Quentin-en-\\[0.08cm]
\HeaderFont\bfseries\itshape Yvelines -- Université Paris-Saclay};

\node[anchor=center] at ([xshift=-2.75cm,yshift=-2.80cm]current page.north east)
{\LogoUVSQ};

\draw[line width=0.75pt]
([xshift=2.80cm,yshift=-5.98cm]current page.north west) --
([xshift=-2.80cm,yshift=-5.98cm]current page.north east);

\node[anchor=north,align=center] at ([yshift=-6.78cm]current page.north)
{\MainTitleFont\bfseries Méthodes de ranking\\[0.15cm]
\MainTitleFont\bfseries et recommandation};

\node[anchor=north,align=center] at ([yshift=-8.65cm]current page.north)
{\SubTitleFont\bfseries Projet 7 : Simulation d'un Google Bombing};

\node[anchor=north,align=center] at ([yshift=-9.45cm]current page.north)
{\FormationFont Rapport de projet};

\draw[line width=0.75pt]
([xshift=2.80cm,yshift=-10.35cm]current page.north west) --
([xshift=-2.80cm,yshift=-10.35cm]current page.north east);

\node[anchor=north,align=center] at ([yshift=-12.10cm]current page.north)
{\BlockTitleFont\bfseries Formation :};

\node[anchor=north,align=center] at ([yshift=-13.05cm]current page.north)
{\FormationFont\bfseries Master 1 Informatique};

\node[anchor=north west,align=left] at ([xshift=3.15cm,yshift=-16.05cm]current page.north west)
{\AuthorTitleFont\bfseries Réalisé par :};

\node[anchor=north west,align=left] at ([xshift=3.65cm,yshift=-16.80cm]current page.north west)
{\AuthorFont \textemdash\; Douadjia Abdelkarim\\
\textemdash\; Issaadi Dorsaf};

\node[anchor=north west,align=left] at ([xshift=11.00cm,yshift=-16.05cm]current page.north west)
{\ResponsibleTitleFont\bfseries Responsable de matière :};

\node[anchor=north west,align=left] at ([xshift=11.50cm,yshift=-16.80cm]current page.north west)
{\ResponsibleFont Monsieur Franck Quessette};

\node[anchor=north,align=center] at ([yshift=-23.30cm]current page.north)
{\FormationFont\bfseries Année universitaire :};

\node[anchor=north,align=center] at ([yshift=-24.35cm]current page.north)
{\YearFont\bfseries 2025 -- 2026};

\end{tikzpicture}
\restoregeometry
\end{titlepage}

\setcounter{page}{1}
\pagestyle{rapportstyle}

\tableofcontents
\newpage

% ============================================================================
\begin{abstract}
Ce rapport présente une simulation de \emph{Google Bombing} dans le cadre du
projet 7. À partir d'un graphe du Web, on calcule d'abord le PageRank initial,
puis on choisit trois cibles : une page de forte pertinence, une page de
pertinence médiane et une page de faible pertinence. On ajoute ensuite trois
structures d'attaque : des sommets isolés, un graphe complet et un anneau.

Les expériences sont menées sur trois graphes de tailles différentes :
\texttt{courtois.txt}, \texttt{Harvard500} et un graphe Wikipedia. Les résultats
montrent que l'attaque par sommets isolés est la plus efficace lorsque le nombre
d'attaquants augmente. L'anneau peut être intéressant pour de petites valeurs de
$k$, tandis que le graphe complet est souvent limité par la dilution interne de
la masse PageRank.
\end{abstract}

% ============================================================================
\section{Introduction}
% ============================================================================

Le \emph{Google Bombing} est une technique de manipulation du référencement qui
exploite la structure des hyperliens du Web pour augmenter artificiellement le
score PageRank d'une page cible. En ajoutant dans le graphe du Web un
sous-graphe composé de pages contrôlées, appelées \emph{attaquants}, pointant
vers la cible, un acteur malveillant peut faire remonter artificiellement la
pertinence d'une page dans un classement basé sur PageRank.

Dans ce projet, nous reprenons l'algorithme PageRank vu en TD et nous le
modifions pour simuler plusieurs structures d'attaque. Le but est de comparer le
PageRank de la cible avant et après l'ajout des attaquants.

Nous faisons varier :
\begin{itemize}[noitemsep]
    \item le nombre $k$ d'attaquants ;
    \item la structure du sous-graphe attaquant : isolés, complet ou anneau ;
    \item le type de cible : forte, médiane ou faible ;
    \item le graphe de départ.
\end{itemize}

Une attaque est considérée comme efficace lorsque la probabilité PageRank de la
cible devient significativement plus forte.

% ============================================================================
\section{Organisation des fichiers du projet}
% ============================================================================

La table~\ref{tab:files} résume les fichiers principaux utilisés dans le projet.

\begin{table}[H]
\centering
\caption{Rôle des fichiers principaux du projet.}
\label{tab:files}
\begin{tabular}{p{0.30\linewidth}p{0.60\linewidth}}
\toprule
\textbf{Fichier ou dossier} & \textbf{Rôle} \\
\midrule
\texttt{googlebombing.c} &
Programme principal en C. Il lit les graphes, calcule le PageRank initial,
ajoute les structures d'attaque et exporte les résultats. \\
\texttt{data/} &
Dossier contenant les graphes utilisés dans les expériences. \\
\texttt{results/} &
Dossier contenant les résultats numériques produits par les simulations. \\
\texttt{images/} &
Dossier contenant les figures utilisées dans le rapport. \\
\texttt{rapport.pdf} &
Rapport final généré à partir du code \LaTeX. \\
\bottomrule
\end{tabular}
\end{table}

% ============================================================================
\section{Rappels théoriques}
% ============================================================================

\subsection{Matrice Google et PageRank}

Soit $\mathcal{G}=(V,E)$ un graphe orienté à $n$ sommets. Chaque sommet
représente une page Web et chaque arc représente un lien hypertexte.

On note $\Hmat \in \mathbb{R}^{n\times n}$ la matrice hyperlien :
\[
H_{ij} =
\begin{cases}
\dfrac{1}{d_i} & \text{si } (i,j)\in E \text{ et } d_i>0,\\[6pt]
0 & \text{sinon,}
\end{cases}
\]
où $d_i$ est le nombre de liens sortants de la page $i$.

\begin{definition}[Matrice Google]
La matrice Google est définie par :
\[
  \G = \alpha \Hmat' + (1-\alpha)\mathbf{e}\mathbf{z}^{\top},
\]
où $\Hmat'$ est la matrice hyperlien corrigée pour les pages sans lien sortant,
$\mathbf{z}$ est le vecteur de personnalisation uniforme, et
$\alpha\in(0,1)$ est le facteur d'amortissement.
\end{definition}

\begin{definition}[Vecteur PageRank]
Le vecteur PageRank $\vpi^*$ est la distribution stationnaire vérifiant :
\[
  \vpi^* = \vpi^*\G,
  \qquad
  \|\vpi^*\|_1=1,
  \qquad
  \vpi^*\geq 0.
\]
\end{definition}

Il est calculé par la méthode des puissances :
\[
  \vpi^{(k+1)}=\vpi^{(k)}\G,
\]
jusqu'à ce que :
\[
  \|\vpi^{(k+1)}-\vpi^{(k)}\|_1 \leq \varepsilon.
\]

\subsection{Convergence sur Harvard500}

Avec $\alpha=0{,}85$ et $\varepsilon=10^{-6}$, la convergence du PageRank initial
sur Harvard500 est atteinte en \textbf{47 itérations}. La norme $L_1$ décroît
de manière géométrique, ce qui correspond au comportement attendu de la méthode
des puissances.

\begin{figure}[H]
\centering
\begin{tikzpicture}
\begin{semilogyaxis}[
  width=0.80\textwidth, height=6cm,
  xlabel={Itération $k$},
  ylabel={Norme $L_1$},
  xmin=0, xmax=46,
  ymin=5e-7, ymax=1.5,
  grid=major, grid style={gray!30},
  title={Convergence du PageRank sur Harvard500 ($\alpha=0{,}85$)},
  title style={font=\small},
  tick label style={font=\small},
  label style={font=\small},
  legend pos=north east, legend style={font=\small},
]
\addplot[blue!70!black, thick] coordinates {
  (0,0.7146)(1,0.4494)(2,0.2959)(3,0.1530)(4,0.0775)
  (5,0.0401)(6,0.0203)(7,0.0110)(8,0.00617)(9,0.00394)
  (10,0.00268)(11,0.00201)(12,0.00153)(13,0.00120)(14,0.000940)
  (15,0.000745)(16,0.000588)(17,0.000467)(18,0.000370)(19,0.000294)
  (20,0.000233)(21,0.000185)(22,0.000147)(23,0.000117)(24,9.35e-5)
  (25,7.46e-5)(26,5.95e-5)(27,4.75e-5)(28,3.79e-5)(29,3.03e-5)
  (30,2.43e-5)(31,1.94e-5)(32,1.56e-5)(33,1.25e-5)(34,1.00e-5)
  (35,8.04e-6)(36,6.46e-6)(37,5.20e-6)(38,4.18e-6)(39,3.37e-6)
  (40,2.72e-6)(41,2.19e-6)(42,1.77e-6)(43,1.43e-6)(44,1.15e-6)
  (45,9.34e-7)
};
\addplot[red, dashed, thin] coordinates {(0,1e-6)(46,1e-6)};
\legend{Norme $L_1$, Seuil $\varepsilon=10^{-6}$}
\end{semilogyaxis}
\end{tikzpicture}
\caption{Décroissance de la norme $L_1$ lors du calcul du PageRank initial sur
Harvard500.}
\label{fig:convergence}
\end{figure}

% ============================================================================
\section{Protocole expérimental}
% ============================================================================

\subsection{Graphes utilisés}

Trois graphes de tailles différentes sont utilisés afin d'évaluer la robustesse
des résultats. Le graphe \texttt{courtois.txt} sert de cas jouet, Harvard500 est
le graphe principal, et Wikipedia permet de tester le passage à grande échelle.

\begin{table}[H]
\centering
\caption{Caractéristiques des graphes expérimentaux.}
\label{tab:graphes}
\begin{tabular}{llccc}
\toprule
\textbf{Graphe} & \textbf{Rôle} & $n$ & $m$ & \textbf{Convergence} \\
\midrule
\texttt{courtois.txt} & Graphe jouet & 8 & 41 & 42 itérations \\
\texttt{Harvard500} & Graphe principal & 500 & 2\,636 & 47 itérations \\
\texttt{Wikipedia} & Grande échelle & 1\,634\,989 & 19\,753\,078 & 35 itérations \\
\bottomrule
\end{tabular}
\end{table}

\subsection{Sélection des cibles}

Après le calcul du PageRank initial, les pages sont triées par score décroissant.
On choisit ensuite trois cibles :
\begin{itemize}[noitemsep]
    \item une cible forte, située dans le haut du classement ;
    \item une cible médiane, située autour du milieu du classement ;
    \item une cible faible, située dans le bas du classement.
\end{itemize}

Pour éviter les cas trop particuliers, les cibles sont choisies par percentile.

\begin{table}[H]
\centering
\caption{Cibles sélectionnées pour chaque graphe.}
\label{tab:cibles}
\small
\begin{tabular}{llccc}
\toprule
\textbf{Graphe} & \textbf{Type} & \textbf{Page} & \textbf{$\PR_0$} & \textbf{Critère} \\
\midrule
\multirow{3}{*}{\texttt{courtois.txt}}
 & Forte   & 6  & $1{,}849 \times 10^{-1}$ & Top 10\% \\
 & Médiane & 4  & $1{,}092 \times 10^{-1}$ & Médiane \\
 & Faible  & 2  & $7{,}641 \times 10^{-2}$ & Bottom 10\% \\
\midrule
\multirow{3}{*}{\texttt{Harvard500}}
 & Forte   & 219 & $2{,}835 \times 10^{-3}$ & Top 10\% \\
 & Médiane & 13  & $6{,}461 \times 10^{-4}$ & Médiane \\
 & Faible  & 121 & $3{,}000 \times 10^{-4}$ & Bottom 10\% \\
\midrule
\multirow{3}{*}{\texttt{Wikipedia}}
 & Forte   & 397700  & $7{,}768 \times 10^{-7}$ & Top 10\% \\
 & Médiane & 694255  & $1{,}429 \times 10^{-7}$ & Médiane \\
 & Faible  & 1471490 & $1{,}009 \times 10^{-7}$ & Bottom 10\% \\
\bottomrule
\end{tabular}
\end{table}

\subsection{Structures d'attaque}

Trois structures de graphes attaquants sont comparées. Dans chaque cas, $k$
désigne le nombre de nouveaux sommets attaquants ajoutés, et la cible est un
sommet déjà présent dans le graphe initial.

\begin{figure}[H]
\centering
\begin{tikzpicture}[
  target/.style={circle, draw=red!70!black, fill=red!12,
                 minimum size=1.1cm, font=\small\bfseries, thick},
  att/.style={circle, draw=blue!65, fill=blue!8,
              minimum size=0.9cm, font=\small},
  >=Stealth, thick
]
\begin{scope}[xshift=0cm]
  \node[target] (T1) at (0,0) {cible};
  \node[att] (A1) at (-1.7,1.7) {$a_1$};
  \node[att] (A2) at (0,2.1)   {$a_2$};
  \node[att] (A3) at (1.7,1.7) {$a_3$};
  \draw[->] (A1)--(T1);
  \draw[->] (A2)--(T1);
  \draw[->] (A3)--(T1);
  \node[below=2pt of T1, font=\footnotesize\itshape] {(a) Isolés};
\end{scope}

\begin{scope}[xshift=5.5cm]
  \node[target] (T2) at (0,0) {cible};
  \node[att] (B1) at (-1.7,1.7) {$a_1$};
  \node[att] (B2) at (0,2.3)   {$a_2$};
  \node[att] (B3) at (1.7,1.7) {$a_3$};
  \draw[->] (B1)--(T2);
  \draw[->] (B2)--(T2);
  \draw[->] (B3)--(T2);
  \draw[->, bend left=18] (B1) to (B2);
  \draw[->, bend left=18] (B2) to (B1);
  \draw[->, bend right=12] (B1) to (B3);
  \draw[->, bend right=12] (B3) to (B1);
  \draw[->, bend left=18] (B2) to (B3);
  \draw[->, bend left=18] (B3) to (B2);
  \node[below=2pt of T2, font=\footnotesize\itshape] {(b) Complet};
\end{scope}

\begin{scope}[xshift=11cm]
  \node[target] (T3) at (0,0) {cible};
  \node[att] (C1) at (-1.7,1.7) {$a_1$};
  \node[att] (C2) at (0,2.5)   {$a_2$};
  \node[att] (C3) at (1.7,1.7) {$a_3$};
  \draw[->] (T3)--(C1);
  \draw[->] (C1)--(C2);
  \draw[->] (C2)--(C3);
  \draw[->] (C3)--(T3);
  \node[below=2pt of T3, font=\footnotesize\itshape] {(c) Anneau};
\end{scope}
\end{tikzpicture}
\caption{Les trois structures de Google Bombing étudiées pour $k=3$ attaquants.}
\label{fig:structures}
\end{figure}

\paragraph{Sommets isolés.}
Chaque attaquant possède un unique lien sortant vers la cible. Toute sa masse
PageRank est donc transmise directement à la cible.

\paragraph{Graphe complet.}
Les attaquants pointent vers la cible, mais aussi entre eux. La masse PageRank
est donc partagée entre la cible et les autres attaquants.

\paragraph{Anneau.}
La cible fait partie d'un cycle avec les attaquants :
\[
a_1 \to a_2 \to \cdots \to a_k \to t \to a_1.
\]
La masse circule en boucle et repasse régulièrement par la cible.

% ============================================================================
\section{Résultats sur Harvard500}
% ============================================================================

\subsection{Évolution du score PageRank absolu}

La figure~\ref{fig:scores-absolus} montre le score PageRank de la cible en
fonction de $k$ pour les trois types de cibles.

\begin{figure}[H]
\centering
\begin{subfigure}[t]{0.48\textwidth}
\begin{tikzpicture}
\begin{axis}[
  width=\textwidth, height=5.5cm,
  xlabel={$k$}, ylabel={Score $\PR$},
  xmin=0, xmax=100,
  grid=major, grid style={gray!25},
  title={\small Cible forte ($\PR_0=0{,}00284$)},
  tick label style={font=\tiny},
  label style={font=\small},
  legend pos=north west,
  legend style={font=\tiny}
]
\addplot[red!80!black, thick] coordinates {
(1,0.0031037)(5,0.0041672)(10,0.0054730)(15,0.0067535)(20,0.0080093)
(25,0.0092412)(30,0.0104499)(40,0.0128001)(50,0.0150649)(60,0.0172488)
(70,0.0193560)(80,0.0213906)(90,0.0233562)(100,0.0252563)};
\addplot[blue!70, thick, dashed] coordinates {
(1,0.0031037)(5,0.0036572)(10,0.0039258)(15,0.0040433)(20,0.0040985)
(25,0.0041222)(30,0.0041281)(40,0.0041106)(50,0.0040730)(60,0.0040256)
(70,0.0039733)(80,0.0039186)(90,0.0038629)(100,0.0038071)};
\addplot[green!60!black, thick, dotted] coordinates {
(1,0.0040502)(5,0.0042859)(10,0.0043717)(15,0.0043793)(20,0.0043581)
(25,0.0043257)(30,0.0042888)(40,0.0042118)(50,0.0041356)(60,0.0040620)
(70,0.0039907)(80,0.0039219)(90,0.0038555)(100,0.0037912)};
\addplot[gray, densely dashed, thin] coordinates {(0,0.002835)(100,0.002835)};
\legend{Isolés, Complet, Anneau, $\PR_0$}
\end{axis}
\end{tikzpicture}
\end{subfigure}
\hfill
\begin{subfigure}[t]{0.48\textwidth}
\begin{tikzpicture}
\begin{axis}[
  width=\textwidth, height=5.5cm,
  xlabel={$k$}, ylabel={Score $\PR$},
  xmin=0, xmax=100,
  grid=major, grid style={gray!25},
  title={\small Cible médiane ($\PR_0=0{,}000646$)},
  tick label style={font=\tiny},
  label style={font=\small},
  legend pos=north west,
  legend style={font=\tiny}
]
\addplot[red!80!black, thick] coordinates {
(1,0.0008998)(5,0.0019049)(10,0.0031390)(15,0.0043491)(20,0.0055360)
(25,0.0067002)(30,0.0078425)(40,0.0100636)(50,0.0122040)(60,0.0142679)
(70,0.0162594)(80,0.0181822)(90,0.0200399)(100,0.0218356)};
\addplot[blue!70, thick, dashed] coordinates {
(1,0.0008998)(5,0.0014304)(10,0.0016996)(15,0.0018279)(20,0.0018978)
(25,0.0019381)(30,0.0019615)(40,0.0019801)(50,0.0019786)(60,0.0019668)
(70,0.0019494)(80,0.0019286)(90,0.0019059)(100,0.0018821)};
\addplot[green!60!black, thick, dotted] coordinates {
(1,0.0011850)(5,0.0018041)(10,0.0020902)(15,0.0021896)(20,0.0022188)
(25,0.0022193)(30,0.0022077)(40,0.0021727)(50,0.0021345)(60,0.0020965)
(70,0.0020598)(80,0.0020242)(90,0.0019899)(100,0.0019568)};
\addplot[gray, densely dashed, thin] coordinates {(0,0.000646)(100,0.000646)};
\legend{Isolés, Complet, Anneau, $\PR_0$}
\end{axis}
\end{tikzpicture}
\end{subfigure}

\vspace{0.5cm}

\begin{subfigure}[t]{0.48\textwidth}
\centering
\begin{tikzpicture}
\begin{axis}[
  width=\textwidth, height=5.5cm,
  xlabel={$k$}, ylabel={Score $\PR$},
  xmin=0, xmax=100,
  grid=major, grid style={gray!25},
  title={\small Cible faible ($\PR_0=0{,}000300$)},
  tick label style={font=\tiny},
  label style={font=\small},
  legend pos=north west,
  legend style={font=\tiny}
]
\addplot[red!80!black, thick] coordinates {
(1,0.0005539)(5,0.0015594)(10,0.0027941)(15,0.0040049)(20,0.0051923)
(25,0.0063571)(30,0.0075000)(40,0.0097222)(50,0.0118636)(60,0.0139286)
(70,0.0159211)(80,0.0178448)(90,0.0197034)(100,0.0215000)};
\addplot[blue!70, thick, dashed] coordinates {
(1,0.0005539)(5,0.0010860)(10,0.0013580)(15,0.0014892)(20,0.0015622)
(25,0.0016056)(30,0.0016320)(40,0.0016566)(50,0.0016609)(60,0.0016548)
(70,0.0016427)(80,0.0016272)(90,0.0016096)(100,0.0015907)};
\addplot[green!60!black, thick, dotted] coordinates {
(1,0.0008672)(5,0.0015200)(10,0.0017818)(15,0.0018669)(20,0.0018909)
(25,0.0018908)(30,0.0018807)(40,0.0018507)(50,0.0018182)(60,0.0017857)
(70,0.0017544)(80,0.0017241)(90,0.0016949)(100,0.0016667)};
\addplot[gray, densely dashed, thin] coordinates {(0,0.000300)(100,0.000300)};
\legend{Isolés, Complet, Anneau, $\PR_0$}
\end{axis}
\end{tikzpicture}
\end{subfigure}

\caption{Score PageRank absolu de la cible en fonction de $k$ sur Harvard500.}
\label{fig:scores-absolus}
\end{figure}

\subsection{Gain relatif}

Pour comparer les attaques indépendamment du score initial, on utilise le gain
relatif :
\[
  \gamma(k)=\frac{\PR_k(t)}{\PR_0(t)}.
\]

\begin{figure}[H]
\centering
\begin{tikzpicture}
\begin{axis}[
  width=0.88\textwidth, height=8cm,
  xlabel={Nombre d'attaquants $k$},
  ylabel={Gain relatif $\gamma = \PR / \PR_0$},
  xmin=0, xmax=100,
  ymin=0, ymax=78,
  grid=major, grid style={gray!20},
  title={Gain relatif selon la structure et le type de cible},
  title style={font=\small},
  tick label style={font=\small},
  label style={font=\small},
  legend pos=north west,
  legend style={font=\footnotesize, cells={anchor=west}, fill=white,
                fill opacity=0.9, draw opacity=1, text opacity=1},
  legend columns=1,
]
\addplot[red!80!black, very thick, solid] coordinates {
(1,1.095)(5,1.470)(10,1.930)(15,2.382)(20,2.825)(25,3.259)
(30,3.686)(40,4.515)(50,5.314)(60,6.084)(70,6.827)(80,7.545)(90,8.238)(100,8.908)};
\addlegendentry{Forte -- Isolés}

\addplot[blue!70!black, very thick, solid] coordinates {
(1,1.393)(5,2.948)(10,4.859)(15,6.732)(20,8.569)(25,10.371)
(30,12.139)(40,15.576)(50,18.889)(60,22.084)(70,25.166)(80,28.142)(90,31.018)(100,33.797)};
\addlegendentry{Médiane -- Isolés}

\addplot[green!60!black, very thick, solid] coordinates {
(1,1.846)(5,5.198)(10,9.314)(15,13.350)(20,17.308)(25,21.191)
(30,25.000)(40,32.407)(50,39.546)(60,46.429)(70,53.070)(80,59.483)(90,65.678)(100,71.667)};
\addlegendentry{Faible -- Isolés}

\addplot[red!80!black, thick, dashed] coordinates {
(1,1.095)(5,1.290)(10,1.385)(15,1.426)(20,1.446)(25,1.454)
(30,1.456)(40,1.450)(50,1.437)(60,1.420)(70,1.401)(80,1.382)(90,1.363)(100,1.343)};
\addlegendentry{Forte -- Complet}

\addplot[blue!70!black, thick, dashed] coordinates {
(1,1.393)(5,2.214)(10,2.631)(15,2.829)(20,2.937)(25,2.999)
(30,3.036)(40,3.065)(50,3.062)(60,3.044)(70,3.017)(80,2.985)(90,2.950)(100,2.913)};
\addlegendentry{Médiane -- Complet}

\addplot[green!60!black, thick, dashed] coordinates {
(1,1.846)(5,3.620)(10,4.527)(15,4.964)(20,5.207)(25,5.352)
(30,5.440)(40,5.522)(50,5.536)(60,5.516)(70,5.476)(80,5.424)(90,5.365)(100,5.302)};
\addlegendentry{Faible -- Complet}

\addplot[red!80!black, thick, dotted] coordinates {
(1,1.429)(5,1.512)(10,1.542)(15,1.545)(20,1.537)(25,1.526)
(30,1.513)(40,1.486)(50,1.459)(60,1.433)(70,1.408)(80,1.383)(90,1.360)(100,1.337)};
\addlegendentry{Forte -- Anneau}

\addplot[blue!70!black, thick, dotted] coordinates {
(1,1.834)(5,2.792)(10,3.235)(15,3.389)(20,3.434)(25,3.435)
(30,3.417)(40,3.363)(50,3.304)(60,3.245)(70,3.188)(80,3.133)(90,3.080)(100,3.029)};
\addlegendentry{Médiane -- Anneau}

\addplot[green!60!black, thick, dotted] coordinates {
(1,2.891)(5,5.067)(10,5.939)(15,6.223)(20,6.303)(25,6.303)
(30,6.269)(40,6.169)(50,6.061)(60,5.952)(70,5.848)(80,5.747)(90,5.650)(100,5.556)};
\addlegendentry{Faible -- Anneau}
\end{axis}
\end{tikzpicture}
\caption{Gain relatif $\gamma$ pour chaque combinaison structure/cible sur
Harvard500.}
\label{fig:gains}
\end{figure}

\subsection{Tableau de synthèse}

\begin{table}[H]
\centering
\caption{Gain relatif $\gamma=\PR/\PR_0$ pour les valeurs clés de $k$ sur Harvard500.}
\label{tab:synthese-harvard}
\small
\begin{tabular}{llrrrr}
\toprule
\textbf{Cible} & \textbf{Structure} & $k=1$ & $k=10$ & $k=50$ & $k=100$ \\
\midrule
\multirow{3}{*}{Forte}
 & Isolés  & $1{,}09\times$ & $1{,}93\times$ & $5{,}31\times$ & $8{,}91\times$ \\
 & Complet & $1{,}09\times$ & $1{,}38\times$ & $1{,}44\times$ & $1{,}34\times$ \\
 & Anneau  & $1{,}43\times$ & $1{,}54\times$ & $1{,}46\times$ & $1{,}34\times$ \\
\midrule
\multirow{3}{*}{Médiane}
 & Isolés  & $1{,}39\times$ & $4{,}86\times$ & $18{,}9\times$ & $33{,}8\times$ \\
 & Complet & $1{,}39\times$ & $2{,}63\times$ & $3{,}06\times$ & $2{,}91\times$ \\
 & Anneau  & $1{,}83\times$ & $3{,}24\times$ & $3{,}30\times$ & $3{,}03\times$ \\
\midrule
\multirow{3}{*}{Faible}
 & Isolés  & $1{,}85\times$ & $9{,}31\times$ & $39{,}5\times$ & $\mathbf{71{,}7}\times$ \\
 & Complet & $1{,}85\times$ & $4{,}53\times$ & $5{,}54\times$ & $5{,}30\times$ \\
 & Anneau  & $2{,}89\times$ & $5{,}94\times$ & $6{,}06\times$ & $5{,}56\times$ \\
\bottomrule
\end{tabular}
\end{table}

% ============================================================================
\section{Analyse et interprétation sur Harvard500}
% ============================================================================

\subsection{Sommets isolés : attaque la plus efficace}

L'attaque par sommets isolés est la plus redoutable. Chaque attaquant possède un
seul lien sortant, dirigé vers la cible. Il n'y a donc pas de perte de masse
vers d'autres attaquants.

La croissance du gain relatif est presque linéaire. Pour la cible faible, le
gain atteint $71{,}7\times$ avec $k=100$. Cela signifie que le score PageRank de
la cible faible est multiplié par plus de 71.

\begin{remarque}
Les sommets isolés paraissent simples, mais ils sont très efficaces parce que
toute la masse PageRank disponible est concentrée vers la cible.
\end{remarque}

\subsection{Graphe complet : saturation par dilution interne}

Le graphe complet contient beaucoup d'arcs, mais il n'est pas optimal. Chaque
attaquant pointe vers la cible, mais aussi vers tous les autres attaquants.
Ainsi, une partie importante de la masse reste dans le sous-graphe attaquant.

Quand $k$ augmente, le nombre d'attaquants augmente, mais la fraction envoyée
par chaque attaquant vers la cible diminue. Cela explique la saturation du gain.

\[
\text{Flux vers la cible}
\simeq
k \times \frac{\PR(a_i)}{k}.
\]

Cette approximation explique pourquoi le graphe complet plafonne rapidement.

\subsection{Anneau : efficace pour de petits cycles}

L'anneau force la masse à revenir régulièrement vers la cible. Pour de petites
valeurs de $k$, le cycle est court et l'effet peut être fort. C'est pour cela
que l'anneau dépasse souvent le graphe complet au début.

Cependant, quand $k$ devient grand, le cycle devient plus long. La masse met
plus de temps à revenir vers la cible, ce qui explique la saturation puis le
léger déclin observé.

\subsection{Effet du type de cible}

Le type de cible joue un rôle central. Le gain relatif s'écrit :
\[
\gamma = 1 + \frac{\Delta \pi}{\PR_0}.
\]

Pour une même augmentation absolue $\Delta \pi$, le gain relatif est plus grand
lorsque $\PR_0$ est petit. C'est la raison pour laquelle les cibles faibles sont
beaucoup plus faciles à booster en ratio que les cibles fortes.

Avec $k=100$ isolés :
\begin{itemize}[noitemsep]
    \item cible forte : gain de $8{,}91\times$ ;
    \item cible médiane : gain de $33{,}8\times$ ;
    \item cible faible : gain de $71{,}7\times$.
\end{itemize}

% ============================================================================
\section{Comparaison multi-graphes}
% ============================================================================

\subsection{Comportement sur \texttt{courtois.txt}}

Sur le petit graphe \texttt{courtois.txt}, les scores PageRank initiaux sont
beaucoup plus élevés que sur Harvard500. L'effet de levier est donc plus faible :
une augmentation absolue de PageRank produit un ratio moins spectaculaire.

\begin{table}[H]
\centering
\caption{Gain relatif $\gamma$ sur \texttt{courtois.txt}.}
\label{tab:courtois}
\small
\begin{tabular}{llrrrr}
\toprule
\textbf{Cible} & \textbf{Structure} & $k=1$ & $k=10$ & $k=50$ & $k=100$ \\
\midrule
\multirow{3}{*}{Forte}
 & Isolés  & $1{,}24\times$ & $2{,}18\times$ & $2{,}83\times$ & $2{,}97\times$ \\
 & Complet & $1{,}24\times$ & $1{,}18\times$ & $0{,}87\times$ & $0{,}26\times$ \\
 & Anneau  & $0{,}43\times$ & $0{,}43\times$ & $0{,}15\times$ & $0{,}08\times$ \\
\midrule
\multirow{3}{*}{Médiane}
 & Isolés  & $1{,}35\times$ & $2{,}75\times$ & $3{,}71\times$ & $3{,}91\times$ \\
 & Complet & $1{,}35\times$ & $1{,}42\times$ & $1{,}05\times$ & $0{,}32\times$ \\
 & Anneau  & $0{,}59\times$ & $0{,}64\times$ & $0{,}23\times$ & $0{,}12\times$ \\
\midrule
\multirow{3}{*}{Faible}
 & Isolés  & $1{,}26\times$ & $2{,}31\times$ & $3{,}04\times$ & $3{,}19\times$ \\
 & Complet & $1{,}26\times$ & $1{,}24\times$ & $0{,}91\times$ & $0{,}27\times$ \\
 & Anneau  & $1{,}06\times$ & $1{,}04\times$ & $0{,}36\times$ & $0{,}20\times$ \\
\bottomrule
\end{tabular}
\end{table}

Un phénomène important apparaît : sur ce très petit graphe, l'anneau peut
diminuer le PageRank de la cible. Cela vient du fait que l'anneau ajoute un lien
sortant depuis la cible vers le premier attaquant. Sur un petit graphe, ce lien
supplémentaire modifie fortement la répartition de la masse.

\subsection{Comportement sur Wikipedia}

Sur le graphe Wikipedia, les expériences sont limitées à $k_{\max}=20$ en raison
du coût de calcul. Malgré cette limite, les tendances restent cohérentes avec
Harvard500 : les sommets isolés donnent le meilleur gain, surtout pour les
cibles faibles.

\begin{table}[H]
\centering
\caption{Gain relatif $\gamma$ sur Wikipedia avec $k_{\max}=20$.}
\label{tab:wikipedia}
\small
\begin{tabular}{llrrrr}
\toprule
\textbf{Cible} & \textbf{Structure} & $k=1$ & $k=5$ & $k=10$ & $k=20$ \\
\midrule
\multirow{3}{*}{Forte}
 & Isolés  & $1{,}11\times$ & $1{,}56\times$ & $2{,}11\times$ & $3{,}22\times$ \\
 & Complet & $1{,}11\times$ & $1{,}35\times$ & $1{,}47\times$ & $1{,}58\times$ \\
 & Anneau  & $1{,}14\times$ & $1{,}43\times$ & $1{,}60\times$ & $1{,}71\times$ \\
\midrule
\multirow{3}{*}{Médiane}
 & Isolés  & $1{,}61\times$ & $4{,}04\times$ & $7{,}08\times$ & $13{,}16\times$ \\
 & Complet & $1{,}61\times$ & $2{,}90\times$ & $3{,}59\times$ & $4{,}16\times$ \\
 & Anneau  & $1{,}83\times$ & $3{,}47\times$ & $4{,}37\times$ & $4{,}91\times$ \\
\midrule
\multirow{3}{*}{Faible}
 & Isolés  & $1{,}85\times$ & $5{,}25\times$ & $9{,}50\times$ & $18{,}00\times$ \\
 & Complet & $1{,}85\times$ & $3{,}66\times$ & $4{,}62\times$ & $5{,}42\times$ \\
 & Anneau  & $1{,}95\times$ & $4{,}27\times$ & $5{,}62\times$ & $6{,}46\times$ \\
\bottomrule
\end{tabular}
\end{table}

\begin{figure}[H]
\centering
\begin{tikzpicture}
\begin{axis}[
  width=0.88\textwidth, height=7cm,
  xlabel={Nombre d'attaquants $k$},
  ylabel={Gain relatif $\gamma$},
  xmin=0, xmax=20,
  ymin=0, ymax=20,
  grid=major, grid style={gray!20},
  title={Gain relatif sur Wikipedia ($n=1{,}6$ M, $\alpha=0{,}85$)},
  title style={font=\small},
  tick label style={font=\small},
  label style={font=\small},
  legend pos=north west,
  legend style={font=\footnotesize, cells={anchor=west}, fill=white,
                fill opacity=0.9, draw opacity=1, text opacity=1},
]
\addplot[red!80!black, very thick, solid] coordinates {
(1,1.11)(2,1.22)(3,1.33)(4,1.44)(5,1.56)(6,1.67)(7,1.78)
(8,1.89)(9,2.00)(10,2.11)(11,2.22)(12,2.33)(13,2.44)(14,2.56)
(15,2.67)(16,2.78)(17,2.89)(18,3.00)(19,3.11)(20,3.22)};
\addlegendentry{Forte -- Isolés}

\addplot[blue!70!black, very thick, solid] coordinates {
(1,1.61)(2,2.22)(3,2.82)(4,3.43)(5,4.04)(6,4.65)(7,5.26)
(8,5.87)(9,6.47)(10,7.08)(11,7.69)(12,8.30)(13,8.91)(14,9.52)
(15,10.12)(16,10.73)(17,11.34)(18,11.95)(19,12.56)(20,13.16)};
\addlegendentry{Médiane -- Isolés}

\addplot[green!60!black, very thick, solid] coordinates {
(1,1.85)(2,2.70)(3,3.55)(4,4.40)(5,5.25)(6,6.10)(7,6.95)
(8,7.80)(9,8.65)(10,9.50)(11,10.35)(12,11.20)(13,12.05)(14,12.90)
(15,13.75)(16,14.60)(17,15.45)(18,16.30)(19,17.15)(20,18.00)};
\addlegendentry{Faible -- Isolés}

\addplot[red!80!black, thick, dashed] coordinates {
(1,1.11)(2,1.19)(3,1.26)(4,1.31)(5,1.35)(6,1.38)(7,1.41)
(8,1.43)(9,1.45)(10,1.47)(11,1.49)(12,1.50)(13,1.52)(14,1.53)
(15,1.54)(16,1.55)(17,1.56)(18,1.56)(19,1.57)(20,1.58)};
\addlegendentry{Forte -- Complet}

\addplot[blue!70!black, thick, dashed] coordinates {
(1,1.61)(2,2.06)(3,2.40)(4,2.68)(5,2.90)(6,3.09)(7,3.24)
(8,3.37)(9,3.49)(10,3.59)(11,3.68)(12,3.75)(13,3.82)(14,3.89)
(15,3.94)(16,3.99)(17,4.04)(18,4.08)(19,4.12)(20,4.16)};
\addlegendentry{Médiane -- Complet}

\addplot[green!60!black, thick, dashed] coordinates {
(1,1.85)(2,2.48)(3,2.96)(4,3.34)(5,3.66)(6,3.91)(7,4.13)
(8,4.32)(9,4.48)(10,4.62)(11,4.74)(12,4.85)(13,4.95)(14,5.03)
(15,5.11)(16,5.19)(17,5.25)(18,5.31)(19,5.37)(20,5.42)};
\addlegendentry{Faible -- Complet}

\addplot[green!60!black, thick, dotted] coordinates {
(1,1.95)(2,2.69)(3,3.31)(4,3.83)(5,4.27)(6,4.64)(7,4.95)
(8,5.21)(9,5.43)(10,5.62)(11,5.78)(12,5.91)(13,6.03)(14,6.12)
(15,6.20)(16,6.27)(17,6.33)(18,6.38)(19,6.43)(20,6.46)};
\addlegendentry{Faible -- Anneau}
\end{axis}
\end{tikzpicture}
\caption{Gain relatif sur Wikipedia. Les sommets isolés progressent
linéairement, tandis que le complet et l'anneau saturent rapidement.}
\label{fig:wikipedia-gains}
\end{figure}

\subsection{Comparaison synthétique des trois graphes}

La figure~\ref{fig:comparaison-graphes} compare le gain relatif de la cible
faible pour l'attaque par sommets isolés. Cette attaque est la plus efficace
dans les expériences principales.

\begin{figure}[H]
\centering
\begin{tikzpicture}
\begin{axis}[
  width=0.85\textwidth, height=6.5cm,
  xlabel={Nombre d'attaquants $k$},
  ylabel={Gain relatif $\gamma$},
  xmin=0, xmax=20,
  ymin=0, ymax=22,
  grid=major, grid style={gray!20},
  title={Comparaison multi-graphes -- cible faible, attaque isolés},
  title style={font=\small},
  tick label style={font=\small},
  label style={font=\small},
  legend pos=north west,
  legend style={font=\small, cells={anchor=west}},
]
\addplot[red!80!black, very thick] coordinates {
(1,1.26)(2,1.47)(3,1.64)(5,1.91)(7,2.10)(10,2.31)(15,2.55)(20,2.69)};
\addlegendentry{\texttt{courtois} ($n=8$)}

\addplot[blue!70!black, very thick] coordinates {
(1,1.846)(2,2.689)(3,3.529)(4,4.365)(5,5.198)(7,6.854)(10,9.314)
(13,12.05)(15,13.35)(17,14.94)(20,17.31)};
\addlegendentry{Harvard500 ($n=500$)}

\addplot[green!60!black, very thick] coordinates {
(1,1.85)(2,2.70)(3,3.55)(4,4.40)(5,5.25)(6,6.10)(7,6.95)
(8,7.80)(9,8.65)(10,9.50)(11,10.35)(12,11.20)(13,12.05)(14,12.90)
(15,13.75)(16,14.60)(17,15.45)(18,16.30)(19,17.15)(20,18.00)};
\addlegendentry{Wikipedia ($n=1{,}6$ M)}
\end{axis}
\end{tikzpicture}
\caption{Comparaison du gain relatif de la cible faible pour l'attaque par
sommets isolés.}
\label{fig:comparaison-graphes}
\end{figure}

\begin{table}[H]
\centering
\caption{Gain relatif $\gamma$ à $k=10$ pour les trois graphes.}
\label{tab:comparaison-k10}
\small
\begin{tabular}{llccc}
\toprule
\textbf{Graphe} & \textbf{Structure} & \textbf{Forte} & \textbf{Médiane} & \textbf{Faible} \\
\midrule
\multirow{3}{*}{\texttt{courtois} ($n=8$)}
 & Isolés  & $2{,}18\times$ & $2{,}75\times$ & $2{,}31\times$ \\
 & Complet & $1{,}18\times$ & $1{,}42\times$ & $1{,}24\times$ \\
 & Anneau  & $0{,}43\times$ & $0{,}64\times$ & $1{,}04\times$ \\
\midrule
\multirow{3}{*}{Harvard500 ($n=500$)}
 & Isolés  & $1{,}93\times$ & $4{,}86\times$ & $9{,}31\times$ \\
 & Complet & $1{,}38\times$ & $2{,}63\times$ & $4{,}53\times$ \\
 & Anneau  & $1{,}54\times$ & $3{,}24\times$ & $5{,}94\times$ \\
\midrule
\multirow{3}{*}{Wikipedia ($n=1{,}6$ M)}
 & Isolés  & $2{,}11\times$ & $7{,}08\times$ & $9{,}50\times$ \\
 & Complet & $1{,}47\times$ & $3{,}59\times$ & $4{,}62\times$ \\
 & Anneau  & $1{,}60\times$ & $4{,}37\times$ & $5{,}62\times$ \\
\bottomrule
\end{tabular}
\end{table}

\subsection{Analyse des différences entre graphes}

Trois phénomènes expliquent les différences observées entre les graphes.

\paragraph{Effet de levier.}
Le score initial des cibles est beaucoup plus élevé sur \texttt{courtois.txt}
que sur Harvard500 ou Wikipedia. Le gain relatif est donc plus faible sur
\texttt{courtois.txt}, car :
\[
\gamma = 1 + \frac{\Delta\pi}{\PR_0}.
\]
Si $\PR_0$ est grand, le même gain absolu produit un ratio plus petit.

\paragraph{Dilution avec la taille du graphe.}
Sur un très grand graphe, la masse de téléportation est distribuée sur beaucoup
plus de sommets. L'effet absolu d'un attaquant est donc plus faible. Cependant,
les cibles faibles ont aussi un PageRank initial très faible, ce qui permet de
garder un gain relatif important.

\paragraph{Cas particulier des petits graphes.}
Sur \texttt{courtois.txt}, l'anneau peut être défavorable. L'ajout d'un arc
sortant depuis la cible vers le premier attaquant modifie fortement le degré
sortant de la cible. Cet effet est beaucoup moins visible sur des graphes plus
grands.

% ============================================================================
\section{Règles empiriques}
% ============================================================================

Les expériences permettent de déduire les règles suivantes :

\begin{enumerate}
    \item \textbf{Les sommets isolés sont l'attaque la plus efficace pour grand $k$.}
    Leur masse est transmise directement vers la cible, sans dilution interne.

    \item \textbf{Le graphe complet est rarement optimal.}
    Les liens internes retiennent une partie de la masse PageRank entre les
    attaquants.

    \item \textbf{L'anneau est intéressant pour petit $k$.}
    Le cycle ramène régulièrement la masse vers la cible, mais l'effet diminue
    lorsque le cycle devient long.

    \item \textbf{Les cibles faibles sont les plus vulnérables en ratio.}
    Leur PageRank initial est faible, donc une petite augmentation absolue donne
    un grand gain relatif.

    \item \textbf{Les très petits graphes peuvent avoir des comportements
    particuliers.}
    Sur \texttt{courtois.txt}, l'ajout d'un arc sortant depuis la cible dans
    l'anneau peut dégrader son score.
\end{enumerate}

% ============================================================================
\section{Conclusion}
% ============================================================================

Cette étude montre qu'un Google Bombing peut augmenter significativement le
PageRank d'une page cible. L'efficacité dépend principalement du nombre
d'attaquants, de la structure du graphe attaquant et du PageRank initial de la
cible.

Sur Harvard500, l'attaque par sommets isolés est la plus efficace. Pour la
cible faible, le gain atteint $71{,}7\times$ avec $k=100$. Pour la cible
médiane, il atteint $33{,}8\times$, tandis que la cible forte reste plus
difficile à faire progresser.

Le graphe complet et l'anneau sont moins efficaces à grande échelle. Le graphe
complet dilue la masse entre les attaquants, tandis que l'anneau est surtout
utile pour de petites valeurs de $k$. Ces résultats confirment que les pages
faibles sont les plus sensibles au Google Bombing et justifient les mécanismes
modernes de détection des fermes de liens.

% ============================================================================
\begin{thebibliography}{9}

\bibitem{brin1998}
S.~Brin, L.~Page,
\textit{The Anatomy of a Large-Scale Hypertextual Web Search Engine},
Computer Networks and ISDN Systems, 30(1--7):107--117, 1998.

\bibitem{langville2006}
A.~N.~Langville, C.~D.~Meyer,
\textit{Google's PageRank and Beyond: The Science of Search Engine Rankings},
Princeton University Press, 2006.

\bibitem{bianchini2005}
M.~Bianchini, M.~Gori, F.~Scarselli,
\textit{Inside PageRank},
ACM Transactions on Internet Technology, 5(1):92--128, 2005.

\bibitem{davis2011}
T.~A.~Davis, Y.~Hu,
\textit{The University of Florida Sparse Matrix Collection},
ACM Transactions on Mathematical Software, 38(1):1--25, 2011.

\end{thebibliography}

\end{document}